\begin{exo}[Formes linéaires sur $\mathcal M_n(\R)$ annulées par des crochets de lie de matrices.]
    \label{al1}
    Soit $\varphi : \mathcal M_n(\R)\to \R$ une forme linéaire vérifiant
    \[\forall M,M'\in \mathcal M_n(\R),\ \varphi(MM')=\varphi(M'M)\quad \text{ et }\quad \varphi(I_n)=n\]
    On pose $A=\lbrace MM'-M'M,\ M,M'\in \mathcal M_n(\R)\rbrace$. 
    \begin{enumerate}
        \item Montrer que $\Vect(A)=\lbrace M\in \mathcal M_n(\R)\ |\ \tr(M)=0\rbrace$.
        \item En déduire que $\varphi=\tr$.
    \end{enumerate}
\end{exo}

\begin{solution}
    \begin{enumerate}
        \item $A\subset \Vect(A)$, donc par définition, 
        $\forall i,j,k,l\in\inl{1}{n},\ E_{ij}E_{kl}-E_{kl}E_{ij}=\delta_{jk}E_{il}-\delta_{li}E_{kj}\in A$. 
    
        Il s'ensuit que
        \[
            \forall i,l\in\inl1n,\ 
            \left\lbrace\begin{array}c i\neq l\implies E_{il}\in \Vect(A)\\
                 \forall i\in\inl{2}{n},\ E_{ii}-E_{11}\in \Vect(A)\end{array}\right. 
        \]
        On pose $\mathcal F=(E_{ij})_{i\neq j}\cup(E_{ii}-E_{11})_{i\in\inl2n}$.
        Cette famille est clairement libre et contient $n^2-1$ vecteurs de $\Vect(A)$,
        d'où $\dim\Vect(A)\geq n^2-1$. 
        Or $\Vect(A)\subset \ker\tr$, et $\tr$ étant une forme linéaire non nulle,
        on a $\dim\ker\tr=n^2-1$, donc $\dim\Vect(A)\leq n^2-1$ 
        donc $\dim\Vect(A)=n^2-1=\dim\ker\tr$ puis $\Vect(A)=\ker\tr$.
        \item De même que pour $\tr$, on a $\Vect(A)\subset \ker\varphi$
        donc $\ker\tr\subset \text{Ker}\varphi$, donc il existe $\lambda\in\R$
        tel que $\varphi=\lambda\text{tr}$. Mais $\varphi(I_n)=n=\text{tr}(I_n)$
        donc $n=\lambda n$ puis $\lambda = 1$ (sauf si $n=0$, mais ce cas est trivial).
    \end{enumerate}
\end{solution}

\begin{exo}[Familles libres de matrices de rang $1$ de $\mathcal M_n(\R)$]
    \label{al2}
	Soit $(X_1,\dots,X_p)\in\R^n$ une famille libre de vecteurs de $\R^n$. Montrer que la famille $(X_1{}^tX_1,\dots,X_p{}^tX_p)$ est une famille libre de vecteurs de $\mathcal M_n(\R)$. Étudier la réciproque.
\end{exo}

\begin{solution}
    Pour $i\in\inl1p$, on note $X_i=(x_i^{(1)},\dots,x_i^{(n)})$. On a alors, \[\forall i\in\inl1p,\ X_i{}^tX_i=(x_i^{(1)}X_i|\dots|x_i^{(n)}X_i)\]
    Soit $(\lambda_1,\dots,\lambda_p)\in\R^p$ telle que $\sum_{i=1}^p\lambda_iX_i{}^tX_i=0$.
    Or, $(X_1,\dots,X_p)$ formant une famille libre, aucun des vecteurs la constituant n'est nul d'où
    \[\forall i\in\inl1p,\ \exists l_i\in\inl1p,\ x_i^{(l_i)}\neq 0\]
    Ainsi, pour $i\in\inl1p$, en regardant que la $l_i$ème colonne dans la somme nulle écrite plus haut, 
    on a \[\sum_{j=1}^p\lambda_jx_j^{(l_i)}X_j=0\]
    $(X_1,\dots,X_p)$ étant libre, on a $\forall j\in\inl1p,\ \lambda_jx_j^{(l_i)}=0$, en particulier $\lambda_ix_i^{(l_i)}=0$
    donc $\lambda_i=0$ ($x_i^{l_i}\neq 0$). 
    Finalement, ceci étant vrai pour tout $i\in\inl1p$, on a 
    \[\forall i\in\inl1p,\ \lambda_i=0\] et $(X_1{}^tX_1,\dots,X_p{}^tX_p)$ est libre dans $\mathcal M_n(\R)$
    (il faut bien sûr préciser que la taille des matrices $X_i{}^tX_i$ est de $n\times n$, mais cela est bien clair).
\end{solution}

\begin{exo}[Fonctions multiplicatives de $\mathcal M_n(\K)$ dans $\K$]
    Soit $f:\mathcal M_n(\K)\to \K$ non constante telle que 
    \[
        \forall A,B\in\mathcal M_n(\K),\ f(AB)=f(A)f(B)    
    \]
    Montrer que $f(A)\neq 0\iff A\in\gl_n(\K)$.
\end{exo}
    
\begin{solution}
    On peut remarquer que $f(I_n) = f(I_n^2) = f(I_n)^2$, alors $f(I_n)\in\lbrace0,1\rbrace$.
    Mais nécessairement $f(I_n) = 1$, car sinon pour toute matrice $M$, $f(M) = f(MI_n)= f(M)f(I_n) = 0$
    et $f$ est l'application nulle, ce qui n'est pas car $f$ n'est pas constante.

    De la même façon, $f(0) = f(0^2) = f(0)^2$, donc $f(0)\in\lbrace0,1\rbrace$. 
    Si $f(0) = 1$, alors pour toute matrice $M$, on a $1=f(0)=f(0\cdot M)=f(0)f(M)=f(M)$, 
    on obtient ainsi que $f$ est constante, ce qui n'est pas. 

    En particulier toute matrice nilpotente $N$ est telle que $f(N)=0$, car si $p\geq 1$
    désigne un entier tel que $N^p=0$, alors $0=f(N^p)=f(N)^p$.

    Soit $A$ une matrice de $\mathcal M_n(\K)$.
    \begin{itemize}
        \item[\boxed{\Leftarrow}] Si $A$ est une matrice inversible, alors  
        \[1=f(I_n)=f(AA^{-1})=f(A)f(A^{-1})\]
        \item[\boxed{\Rightarrow}] Supposons que $A$ est telle que $f(A)\neq 0$.
        On note $r$ le rang de $A$, on veut montrer que $r=n$.

        On sait que $A$ est équivalente à $J_r$ : il existe des matrices inversibles $P$ et $Q$
        qui vérifient $A = PJ_r Q$. 
        Alors $f(A) = f(P J_r Q) = f(P)f(J_r)f(Q)$ et donc $f(J_r)\neq0$.
        
        Si $r<n$, alors il existe une matrice nilpotente de rang $r$ (considérer une matrice où on met des coefficients non nuls seulement sur la sur-diagonale)
        aussi équivalente à $J_r$. 

        Comme on a déjà montré que les inversibles ont une image non nulle par $f$, il est clair que $f(J_r)=0$, ce qui n'est pas.
        En conclusion, $r=n$ et $A$ est inversible.
    \end{itemize}
\end{solution}